% Options for packages loaded elsewhere
% Options for packages loaded elsewhere
\PassOptionsToPackage{unicode}{hyperref}
\PassOptionsToPackage{hyphens}{url}
\PassOptionsToPackage{dvipsnames,svgnames,x11names}{xcolor}
%
\documentclass[
  french,
  11pt,
  a4paper,
  twoside,
  article,
  titlepage=false]{memoir}
\usepackage{xcolor}
\usepackage{amsmath,amssymb}
\setcounter{secnumdepth}{-\maxdimen} % remove section numbering
\usepackage{iftex}
\ifPDFTeX
  \usepackage[T1]{fontenc}
  \usepackage[utf8]{inputenc}
  \usepackage{textcomp} % provide euro and other symbols
\else % if luatex or xetex
  \usepackage{unicode-math} % this also loads fontspec
  \defaultfontfeatures{Scale=MatchLowercase}
  \defaultfontfeatures[\rmfamily]{Ligatures=TeX,Scale=1}
\fi
\usepackage{lmodern}
\ifPDFTeX\else
  % xetex/luatex font selection
\fi
% Use upquote if available, for straight quotes in verbatim environments
\IfFileExists{upquote.sty}{\usepackage{upquote}}{}
\IfFileExists{microtype.sty}{% use microtype if available
  \usepackage[]{microtype}
  \UseMicrotypeSet[protrusion]{basicmath} % disable protrusion for tt fonts
}{}
\makeatletter
\@ifundefined{KOMAClassName}{% if non-KOMA class
  \IfFileExists{parskip.sty}{%
    \usepackage{parskip}
  }{% else
    \setlength{\parindent}{0pt}
    \setlength{\parskip}{6pt plus 2pt minus 1pt}}
}{% if KOMA class
  \KOMAoptions{parskip=half}}
\makeatother
% Make \paragraph and \subparagraph free-standing
\makeatletter
\ifx\paragraph\undefined\else
  \let\oldparagraph\paragraph
  \renewcommand{\paragraph}{
    \@ifstar
      \xxxParagraphStar
      \xxxParagraphNoStar
  }
  \newcommand{\xxxParagraphStar}[1]{\oldparagraph*{#1}\mbox{}}
  \newcommand{\xxxParagraphNoStar}[1]{\oldparagraph{#1}\mbox{}}
\fi
\ifx\subparagraph\undefined\else
  \let\oldsubparagraph\subparagraph
  \renewcommand{\subparagraph}{
    \@ifstar
      \xxxSubParagraphStar
      \xxxSubParagraphNoStar
  }
  \newcommand{\xxxSubParagraphStar}[1]{\oldsubparagraph*{#1}\mbox{}}
  \newcommand{\xxxSubParagraphNoStar}[1]{\oldsubparagraph{#1}\mbox{}}
\fi
\makeatother


\usepackage{longtable,booktabs,array}
\usepackage{calc} % for calculating minipage widths
% Correct order of tables after \paragraph or \subparagraph
\usepackage{etoolbox}
\makeatletter
\patchcmd\longtable{\par}{\if@noskipsec\mbox{}\fi\par}{}{}
\makeatother
% Allow footnotes in longtable head/foot
\IfFileExists{footnotehyper.sty}{\usepackage{footnotehyper}}{\usepackage{footnote}}
\makesavenoteenv{longtable}
\usepackage{graphicx}
\makeatletter
\newsavebox\pandoc@box
\newcommand*\pandocbounded[1]{% scales image to fit in text height/width
  \sbox\pandoc@box{#1}%
  \Gscale@div\@tempa{\textheight}{\dimexpr\ht\pandoc@box+\dp\pandoc@box\relax}%
  \Gscale@div\@tempb{\linewidth}{\wd\pandoc@box}%
  \ifdim\@tempb\p@<\@tempa\p@\let\@tempa\@tempb\fi% select the smaller of both
  \ifdim\@tempa\p@<\p@\scalebox{\@tempa}{\usebox\pandoc@box}%
  \else\usebox{\pandoc@box}%
  \fi%
}
% Set default figure placement to htbp
\def\fps@figure{htbp}
\makeatother



\ifLuaTeX
\usepackage[bidi=basic]{babel}
\else
\usepackage[bidi=default]{babel}
\fi
% get rid of language-specific shorthands (see #6817):
\let\LanguageShortHands\languageshorthands
\def\languageshorthands#1{}


\setlength{\emergencystretch}{3em} % prevent overfull lines

\providecommand{\tightlist}{%
  \setlength{\itemsep}{0pt}\setlength{\parskip}{0pt}}



 


\makepagestyle{style-td1}
\makeevenhead{style-td1}{\sffamily\small Intégration et Probabilités}{\sffamily B}{\sffamily\small CC I MIASHS 1 \& 2}
\makeoddhead{style-td1}{\sffamily\small Intégration et Probabilités}{\sffamily B}{\sffamily\small CC I MIASHS 1 \& 2}
\makeevenfoot{style-td1}{\sffamily\small MA1AY010}{\thepage}{\sffamily\small L3 MIASHS}
\makeoddfoot{style-td1}{\sffamily\small L3 MIASHS}{\thepage}{\sffamily\small MA1AY010}
\pagestyle{style-td1}
\makeatletter
\@ifpackageloaded{tcolorbox}{}{\usepackage[skins,breakable]{tcolorbox}}
\@ifpackageloaded{fontawesome5}{}{\usepackage{fontawesome5}}
\definecolor{quarto-callout-color}{HTML}{909090}
\definecolor{quarto-callout-note-color}{HTML}{0758E5}
\definecolor{quarto-callout-important-color}{HTML}{CC1914}
\definecolor{quarto-callout-warning-color}{HTML}{EB9113}
\definecolor{quarto-callout-tip-color}{HTML}{00A047}
\definecolor{quarto-callout-caution-color}{HTML}{FC5300}
\definecolor{quarto-callout-color-frame}{HTML}{acacac}
\definecolor{quarto-callout-note-color-frame}{HTML}{4582ec}
\definecolor{quarto-callout-important-color-frame}{HTML}{d9534f}
\definecolor{quarto-callout-warning-color-frame}{HTML}{f0ad4e}
\definecolor{quarto-callout-tip-color-frame}{HTML}{02b875}
\definecolor{quarto-callout-caution-color-frame}{HTML}{fd7e14}
\makeatother
\makeatletter
\@ifpackageloaded{caption}{}{\usepackage{caption}}
\AtBeginDocument{%
\ifdefined\contentsname
  \renewcommand*\contentsname{Table des matières}
\else
  \newcommand\contentsname{Table des matières}
\fi
\ifdefined\listfigurename
  \renewcommand*\listfigurename{Liste des Figures}
\else
  \newcommand\listfigurename{Liste des Figures}
\fi
\ifdefined\listtablename
  \renewcommand*\listtablename{Liste des Tables}
\else
  \newcommand\listtablename{Liste des Tables}
\fi
\ifdefined\figurename
  \renewcommand*\figurename{Figure}
\else
  \newcommand\figurename{Figure}
\fi
\ifdefined\tablename
  \renewcommand*\tablename{Table}
\else
  \newcommand\tablename{Table}
\fi
}
\@ifpackageloaded{float}{}{\usepackage{float}}
\floatstyle{ruled}
\@ifundefined{c@chapter}{\newfloat{codelisting}{h}{lop}}{\newfloat{codelisting}{h}{lop}[chapter]}
\floatname{codelisting}{Listing}
\newcommand*\listoflistings{\listof{codelisting}{Liste des Listings}}
\usepackage{amsthm}
\theoremstyle{definition}
\newtheorem{exercise}{Exercice}[section]
\theoremstyle{remark}
\AtBeginDocument{\renewcommand*{\proofname}{Preuve}}
\newtheorem*{remark}{Remarque}
\newtheorem*{solution}{Solution}
\newtheorem{refremark}{Remarque}[section]
\newtheorem{refsolution}{Solution}[section]
\makeatother
\makeatletter
\makeatother
\makeatletter
\@ifpackageloaded{caption}{}{\usepackage{caption}}
\@ifpackageloaded{subcaption}{}{\usepackage{subcaption}}
\makeatother
\usepackage{bookmark}
\IfFileExists{xurl.sty}{\usepackage{xurl}}{} % add URL line breaks if available
\urlstyle{same}
\hypersetup{
  pdflang={fr},
  colorlinks=true,
  linkcolor={blue},
  filecolor={Maroon},
  citecolor={Blue},
  urlcolor={Blue},
  pdfcreator={LaTeX via pandoc}}


\author{}
\date{}
\begin{document}
\frontmatter

\counterwithout{exercise}{section}


\mainmatter
\begin{tcolorbox}[enhanced jigsaw, colback=white, left=2mm, arc=.35mm, breakable, rightrule=.15mm, bottomrule=.15mm, toprule=.15mm, colframe=quarto-callout-note-color-frame, opacityback=0, leftrule=.75mm]

\vspace{-3mm}\textbf{CC I/IV :}\vspace{3mm}

\begin{itemize}
\tightlist
\item
  Licence MIASHS
\item
  20 Octobre 2025 Durée : 1 heure 15
\item
  \textbf{Intégration et Probabilités}
\end{itemize}

\end{tcolorbox}

\begin{exercise}[]\protect\hypertarget{exr-1-a}{}\label{exr-1-a}

~

\end{exercise}

\textbf{Partie I}

Dans un étang vivent \(N\geqslant 1\) poissons numérotés \(1\) à \(N\).
On pêche le poisson \(i\) avec une probabilité \(p_i\)
(\(\sum_{i=1}^N  p_i=1\)).\\
Après avoir noté le numéro du poisson pêché, on le remet (délicatement)
à l'eau. On répète cette opération \(n\geq 1\) fois. Les opérations sont
indépendantes les unes des autres.

\begin{enumerate}
\def\labelenumi{\roman{enumi}.}
\tightlist
\item
  Proposer un univers, une tribu et une loi sur la tribu pour modéliser
  ces pêches répétées.
\item
  Quelle est la probabilité que le poisson \(i\) soit pêché au moins une
  fois?
\item
  Quelle est la probabilité que le poisson \(i\) soit pêché exactement
  \(k\) fois?
\item
  Quelle est la probabilité que le poisson pêché lors de l'épreuve \(n\)
  n'ait pas été pêché lors des \(n-1\) épreuves précédentes.
\end{enumerate}

\textbf{Partie II}

Dans le même étang, au lieu de répéter l'épreuve \(n\) fois, on tire
d'abord \(X\) selon une loi de Poisson de paramètre \(n\). Si \(X=k\),
on effectue ensuite \(k\) pêches indépendantes selon la description plus
haut.

\begin{enumerate}
\def\labelenumi{\alph{enumi}.}
\setcounter{enumi}{21}
\tightlist
\item
  Proposer un univers, une tribu et une loi sur la tribu pour modéliser
  ces pêches répétées un nombre aléatoire de fois.
\end{enumerate}

\begin{enumerate}
\def\labelenumi{\roman{enumi}.}
\setcounter{enumi}{5}
\tightlist
\item
  Montrer que pour chaque poisson \(i\), le nombre captures du poisson
  \(i\) suit une loi de Poisson. Préciser le paramètre de la loi de
  Poisson. Quelle est la probabilité que le poisson \(i\) soit pêché
  exactement \(\ell\) fois?
\item
  Quelle est la probabilité que le poisson \(i\) soit pêché exactement
  \(\ell\) fois \emph{et} le poisson \(j\) soit pêché exactement \(m\)
  fois ?
\end{enumerate}

\emph{Rappel} : la loi de Poisson de paramètre \(\mu>0\) est définie par

\[P(\{k\}) =  \mathrm{e}^{-\mu} \frac{\mu^k}{k!} \qquad \text{ pour } k \in \mathbb{N}\, .\]

\begin{exercise}[]\protect\hypertarget{exr-5-a}{}\label{exr-5-a}

~

\end{exercise}

Soit \(\alpha >0\), \(\mu\) la mesure sur
\((\mathbb{R}, \mathcal{B}(\mathbb{R}))\) telle que pour tout \(a <  b\)

\[\mu((a,b)) = \alpha \int_a^b (2-|x|)^2 \times \mathbb{I}_{[-2,2]}(x) \mathrm{d}x\]

avec

\[\mathbb{I}_{[c,d]}(x) = \begin{cases} 1 & \text{ si } x \in [c,d] \\ 0 & \text{sinon}\end{cases}\]

\begin{enumerate}
\def\labelenumi{\roman{enumi}.}
\tightlist
\item
  Calculer \(\mu(0,1)\), \(\mu(1, 3)\)
\item
  Calculer la ou les valeurs de \(\alpha\) pour que \(\mu\) soit une
  probabilité
\item
  Calculer \(\mu([a,b])\) pour \(a<b\).
\end{enumerate}

\newpage{}

\begin{exercise}[]\protect\hypertarget{exr-6-a}{}\label{exr-6-a}

~

\end{exercise}

Dans une urne on a disposé \(3\) boules numérotées de \(1\) à \(3\). On
effectue \(3\) tirages successifs uniformes \emph{sans remises}. Les
tirages sont numérotés de \(1\) à \(3\) (rangs des tirages). À chaque
tirage, on observe le numéro de la boule obtenue.

\begin{enumerate}
\def\labelenumi{\roman{enumi}.}
\tightlist
\item
  Décrire un espace \(\Omega, \mathcal{F}, P\) pour analyser
  l'expérience aléatoire.
\item
  Quelle est la probabilité que la boule \(i\) soit observée au tirage
  \(i\) (rang \(i\)) ?
\item
  Pour \(i \neq j\), quelle est la probabilité que la boule \(i\) soit
  observée au tirage \(i\) et la boule \(j\) au tirage \(j\) ?
\item
  Quelle est la probabilité que pour tous les \(3\) tirages, le numéro
  de la boule observée coïncide avec le rang du tirage ?
\item
  Quelle est la probabilité que pour au moins un des \(3\) tirages, le
  numéro de la boule observée coïncide avec le rang du tirage ?
\item
  Quelle est la probabilité que pour tous les \(3\) tirages, le numéro
  de la boule observée soit différent du rang du tirage ?
\end{enumerate}


\backmatter


\end{document}
